\documentclass[11pt]{article}
\usepackage[T1]{fontenc}
\usepackage{lmodern,microtype,amsmath,amssymb,amsthm}
\usepackage[margin=1in]{geometry}
\usepackage{booktabs,longtable,tabularx,array,listings,xcolor}
\usepackage[hidelinks]{hyperref}
\hypersetup{pdftitle={Accelerating Exact Unknot Recognition},
 pdfauthor={Technical report prepared for Vladimir Reshetnikov},
 pdfsubject={Implementation, correctness arguments, and reproducible benchmarks},
 pdfkeywords={unknot recognition, Khovanov homology, Jones polynomial, optimization}}
\usepackage{fancyhdr}
\setlength{\headheight}{14pt}
\pagestyle{fancy}\fancyhf{}
\fancyhead[L]{\small Accelerating exact unknot recognition}
\fancyhead[R]{\small Technical report}
\fancyfoot[C]{\thepage}
\newtheorem{theorem}{Theorem}[section]
\newtheorem{proposition}[theorem]{Proposition}
\newtheorem{lemma}[theorem]{Lemma}
\newtheorem{corollary}[theorem]{Corollary}
\theoremstyle{remark}\newtheorem{remark}[theorem]{Remark}
\newcommand{\F}{\mathbb F}
\newcommand{\Kh}{\operatorname{Kh}}
\newcommand{\rKh}{\widetilde{\operatorname{Kh}}}
\newcommand{\OO}{\mathcal O}
\newcommand{\code}[1]{\texttt{\small #1}}
\lstset{basicstyle=\ttfamily\small,columns=fullflexible,keepspaces=true,
  breaklines=true,frame=single,framerule=0.3pt,showstringspaces=false,
  aboveskip=10pt,belowskip=10pt}
\setlength{\emergencystretch}{2em}
\setcounter{tocdepth}{1}
\newcommand{\HardTime}{2.23}
\newcommand{\TimeLimit}{30}
\newcommand{\TestCount}{37}
\newcommand{\HardRank}{2949}
\newcommand{\HardUnreduced}{5898}
\newcommand{\HardPeakObjects}{17694}
\newcommand{\HardWidth}{8}
\newcommand{\PythonVersion}{Python 3.13.5}
\newcommand{\HostPlatform}{Linux-6.18.44-x86\_64-with-glibc2.41}
\newcommand{\CpuModel}{INTEL(R) XEON(R) PLATINUM 8573C}
\title{Accelerating Exact Unknot Recognition\\[5pt]
\large Incremental preprocessing, modular obstructions,\\and sparse cobordism elimination}
\author{Technical report prepared for Vladimir Reshetnikov}
\date{18 September 2026}
\begin{document}
\maketitle
\begin{abstract}
We modify the pure-Python recognizer in the supplied \code{Knots.zip/fast/}
rather than substitute an unrelated package. The resulting implementation is
complete when its optional resource limits are disabled, retains the scanning
Khovanov backend over $\F_2$, and has no third-party runtime dependencies.
The changes comprise incremental Reidemeister reductions and scan ordering,
a linear-time descending-diagram test, two sound modular rejection tests,
cached cobordism topology, characteristic-two algebraic shortcuts, and
sparsity-aware Gaussian cancellation. The supplied 36-crossing benchmark now
has reduced Khovanov rank $2949$ computed in \HardTime\ seconds (median of
three cold-cache runs); the unchanged baseline reaches our \TimeLimit-second
limit on the same host. This is a raw homology computation, not merely a
prefilter verdict. All \TestCount\ tests pass, including independent cube
calculations and differential checks. Some tiny inputs become slower, and the
full benchmark tables include them. We prove asymptotic improvements for
several subroutines, but do \emph{not} claim a quasi-polynomial worst-case
bound for the complete recognizer. We explain precisely why the hierarchy
materials do not transfer such a bound to this implementation.
\end{abstract}
\vspace{0.5em}

\section{Deliverable and scope}

The input is a validated, one-component planar knot diagram with $n$ crossings.
The package accepts the original planar-diagram (PD), closed-braid, and grid
JSON formats. The result is \code{UNKNOT}, \code{KNOTTED}, or \code{UNKNOWN}.
The last result means that a requested resource ceiling was reached, not that
an invariant was trivial. With no time or Khovanov-object ceiling, the pipeline
is a decision procedure. Internal Jones-filter ceilings merely abandon that
filter and continue to the complete backend.

The implementation is in \code{fastunknot/}; the unmodified supplied Python
package is preserved in \code{baseline/fastunknot/}. The files
\code{diagram.py} and \code{alexander.py} are unchanged. The new
\code{filters.py} contains the modular algorithms. The main changes to
\code{scan.py}, \code{simplify.py}, and \code{recognize.py} are described below.
Tests, fixed benchmark inputs, an isolated-process comparison driver, raw
results, and the sources for this report are included. The existing MIT
No Attribution license is retained. The archive manifest records file hashes;
provenance metadata identifies the supplied archive and the unchanged sources.

This is not an implementation of Lackenby's hierarchy algorithm. Its strongest
achieved result is a substantial practical improvement to the supplied exact
algorithm, together with genuine better asymptotic bounds for preprocessing
and order construction. The worst-case envelope for the complete implementation
remains $2^{\OO(n)}$. The rank routine computes total rank and cube-degree
ranks, not the full quantum-graded Poincar\'e polynomial or integral torsion.

\section{What the hierarchy materials establish}
\label{sec:lackenby}

The supplied talk's main theorem slide, physical page 14, explicitly announces
an $n^{\OO(\log n)}$ unknot algorithm and labels the result ``L 2021''
\cite{lackenby-talk}. The supplied July 2026 paper presents a hierarchy-based
algorithm for incompressibility and obtains unknot recognition as a special
case \cite{lackenby-paper}. Its introduction describes further speed-up as
future work. Section 9 deliberately gives only a crude iteration estimate;
Section 10 refines the construction using interior parallelity bundles and
controls hierarchy length in terms of the initial complexity. These are
important distinctions: an announcement, a terminating algorithm, and a
particular implementation with a proved running-time bound are different
objects.

In the notation of Section 9, one relevant bound is
\begin{equation}\label{eq:hierarchy}
  \text{number of iterations}\ \leq L(g+1)^L,
\end{equation}
where $L$ bounds hierarchy length and $g$ bounds a surface complexity.
A linear bound on $L$ in an initial input-complexity parameter is not, by
itself, a logarithmic bound. Nor is an iteration count a bit-complexity
bound until representation sizes and the operations inside an iteration
are controlled.

\paragraph{A quantitative implementation target.}
For the \emph{crude bound} \eqref{eq:hierarchy}, the hypothetical conditions
$L\leq a\log n$, $g\leq n^b$, and $n^c$ work per iteration would imply
\[
  n^c L(g+1)^L
  \leq n^c a\log n\,(n^b+1)^{a\log n}
  =\exp\bigl(\OO((\log n)^2)\bigr).
\]
This is a sufficient set of conditions, not a claim that these bounds have
been established for the supplied hierarchy procedure, or that every route
to quasi-polynomial time must use them. A different sufficient contract would
be a recurrence
\[
 T(N)\leq N^a T(\rho N)+N^b,\qquad 0<\rho<1,
\]
with polynomial-size encodings and polynomially bounded branching. Its
$\OO(\log N)$ levels give the same quasi-polynomial envelope.

Turning either contract into code would require more than a normal-surface
search called behind an opaque interface. A useful implementation boundary
would consist of the following concrete data and operations: an incidence
encoding of handle structures and boundary patterns; binary multiplicities
for families of parallel normal pieces; splitting and regluing operations on
these compressed families; recognition and removal of parallelity bundles;
and hierarchical multi-surfaces with explicitly tracked complexity. Each
operation needs an invariant explaining what manifold and boundary pattern
its output represents. The balancing or Cheeger-region step must construct
an admissible surface and prove the required decrease, rather than merely
return a graph separator. Finally, the number of operations and the bit
length of every multiplicity need bounds in the original $n$.

None of those geometric operations is supplied by the Khovanov engine.
A small scan boundary, a small cache, or a successful numerical experiment
cannot fill that gap. This work therefore improves the available engine
instead of labeling an incomplete hierarchy interface a quasi-polynomial
implementation.

\section{Pipeline and its decision invariant}

The optimized pipeline is
\[
 \begin{gathered}
 \text{R1/R2 simplification}\ \longrightarrow\ \text{descending test}\\
 \longrightarrow\ \text{modular determinant}\ \longrightarrow\
 \text{bounded modular bracket}\\
 \longrightarrow\ \text{exact Alexander polynomial}\ \longrightarrow\
 \text{scanning Khovanov rank over }\F_2.
 \end{gathered}
\]
The first two stages may establish triviality. The next three may establish
nontriviality. Only the last stage is a complete invariant test for this
purpose. A modular residue that matches the unknot is never interpreted as
an unknot verdict. In particular, the implementation does not assume that
the Jones polynomial detects the unknot.

\begin{theorem}[Decision rule]
Assume the input satisfies the package's knot-diagram validation conditions
and the implemented cobordism operations represent the stated Khovanov chain
operations. Every completed pipeline verdict is sound. In the absence of
user-imposed time and object limits, the pipeline terminates with a verdict.
\end{theorem}
\begin{proof}
Crossing-decreasing Reidemeister moves preserve the represented knot, and a
descending diagram represents the unknot. Sections~\ref{sec:det} and
\ref{sec:jones} prove the modular implications. The inherited exact Alexander
stage rejects only a normalized polynomial different from $1$.

For the final stage, Shumakovitch's splitting over $\F_2$ implies
\[
 \dim_{\F_2}\Kh(K;\F_2)=2\dim_{\F_2}\rKh(K;\F_2)
\]
\cite[Corollary 3.2.C]{shumakovitch}. Universal coefficients imply that the
rational reduced rank is at most the reduced rank over $\F_2$; the former
is nonzero. Kronheimer--Mrowka's theorem identifies reduced rational rank
one with the unknot \cite{kronheimer-mrowka}. Thus reduced $\F_2$ rank one
implies triviality, while the unknot itself has reduced $\F_2$ rank one.

The finite cube of resolutions and exact finite-dimensional algebra give a
terminating fallback. The bounded Jones stage cannot prevent reaching it.
Every earlier successful stage makes a sound decision, so the composition
has the same property.
\end{proof}

This theorem states the mathematical contract, not a machine-checked proof
of every Python line. The implementation is validated by the independent
and differential tests in Section~\ref{sec:validation}. The JSON evidence
contains the input digest, simplification trace, residues or polynomial,
and, when applicable, scan order and rank data. A homology result is a
reproducible computation, not a short independently verified certificate.

\section{Incremental preprocessing}
\label{sec:preprocessing}

\subsection{A mutable dart involution for R1 and R2}

Give each crossing four permanent dart identifiers $4c,\ldots,4c+3$.
The involution $\alpha$ pairs darts along edges, while
$\sigma(4c+i)=4c+(i+1\bmod4)$ rotates counterclockwise around a crossing.
Faces are orbits of $F=\sigma\alpha$. Therefore a removable monogon is
recognized by $F(d)=d$. A candidate bigon satisfies $F^2(d)=d$ with its
two darts at distinct crossings; the additional parity test checks that
one strand is over at both crossings. Odd ports are over-passes in the
inherited PD convention.

The old implementation searches all faces and rebuilds the complete diagram
after each move. The new one stores legal local candidates in a heap. When
one or two crossings are removed, their opposite ports are joined through
the deleted region. At most eight deleted darts need inspection; only a
constant number of external involution entries change. Every newly created
monogon or bigon must meet that changed neighborhood. Rechecking all four
darts of every affected external crossing therefore finds every newly legal
move. Stale heap entries are discarded on extraction.

Heap keys use move kind followed by permanent crossing identifiers, retaining
the old R1-before-R2 and lexicographic policy. Permanent identifiers retain
relative order after deletion. A Fenwick tree converts them to indices in
the \emph{current} diagram for the original replayable \code{Move} format.
This avoids rebuilding or renumbering the surviving diagram after every
move. A final reconstruction calls the original validator once.

\begin{proposition}
The incremental simplifier produces the same move sequence as the baseline
policy, up to harmless edge-label renaming, and uses $\OO(n\log n)$ word
operations and $\OO(n)$ space.
\end{proposition}
\begin{proof}
The initial candidates are exhaustive. Locality of the face successor map
establishes the candidate invariant after every deletion. Permanent ordering
agrees with current ordering, and the Fenwick prefix sum gives its exact
current rank; hence each chosen valid move is the baseline's first legal
move. There are at most $n$ deletions, each creates a bounded number of
heap entries, and both heap and prefix-sum work cost $\OO(\log n)$ per
operation. The initial and final passes are at most $\OO(n\log n)$.
\end{proof}

Here ``word operations'' treats normalized labels and indices of
$\OO(\log n)$ bits as machine words. The statement does not assert constant
cost for arbitrary-precision integers of unbounded length. The old repeated
full-diagram work has a quadratic lower bound on the family of successive
curls used in the preprocessing benchmark.

\subsection{Descending diagrams by interval coverage}

Fix an oriented traversal with $2n$ visits. For a crossing let $o$ be its
over-visit position and $u$ its under-visit position. The cyclic starting
positions for which this crossing is encountered underneath first form
exactly the interval $(o,u]$. Add all such intervals to a circular difference
array. A zero in the prefix sum is a start for which every crossing is first
encountered over. Repeat for the reverse orientation.

\begin{proposition}
A descending basepoint, when one exists in either traversal direction, is
found in $\OO(n)$ time and space.
\end{proposition}
\begin{proof}
A start is valid exactly when it belongs to none of the forbidden intervals.
Each interval is represented by at most two ordinary intervals, hence a
constant number of difference-array updates. Two prefix scans and two
traversals suffice. The returned dart identifies the actual basepoint; the
algorithm is not merely a boolean relaxation.
\end{proof}

\section{A modular determinant before polynomial arithmetic}
\label{sec:det}

The inherited Alexander implementation computes an exact polynomial using
fraction-free elimination. That is useful but unnecessary for many negative
instances. At $t=-1$, a Wirtinger Fox row has the form
\[
 2x_{\mathrm{over}}-x_{\mathrm{under},1}-x_{\mathrm{under},2}=0.
\]
Over-passing edge segments are joined by disjoint-set union to construct the
Wirtinger arcs. Remove one row and one column of the resulting matrix, and
compute its determinant in $\F_p$, with the fixed prime
$p=1\,000\,000\,007$. Sparse dictionaries represent rows; among candidate
pivots in a column the algorithm picks a row with the fewest nonzero
entries. Row swaps and pivot products are retained, so the computation is
a determinant, not just a rank test.

\begin{proposition}
If the computed residue is neither $1$ nor $-1$ modulo $p$, the input knot
is nontrivial.
\end{proposition}
\begin{proof}
The Fox minor evaluates to $\Delta_K(-1)$ up to sign; multiplication of the
Alexander polynomial by $\pm t^k$ also gives only a sign at $t=-1$. For the
unknot the normalized Alexander polynomial is $1$, so every allowed signed
value is $\pm1$. Reduction modulo $p$ preserves equality. A different
residue therefore contradicts triviality.
\end{proof}

This filter is deterministic and one-sided. A nontrivial determinant can
be congruent to $\pm1$; that is an inconclusive result, not an error. The
worst-case cost is $\OO(n^3)$ field operations and $\OO(n^2)$ stored field
elements. The prime is fixed, so coefficient swell is eliminated at this
stage. The exact polynomial remains available for cases that pass both
modular filters.

\section{A bounded, independent Jones-bracket engine}
\label{sec:jones}

\subsection{Normalization and soundness}

For ports $(a,b,c,d)$ in the inherited counterclockwise PD convention, use
smoothings $0=((a,b),(c,d))$ and $1=((a,d),(b,c))$ with coefficients $A$ and
$A^{-1}$. Set $\delta=-A^2-A^{-2}$. It is convenient to count \emph{every}
closed circle, including the first, giving the all-loops bracket
\begin{equation}\label{eq:bracket}
 B_D(A)=\sum_{s\in\{0,1\}^n} A^{n-2|s|}\delta^{c(s)},
\end{equation}
where $c(s)$ is the number of circles in state $s$ and an empty PD denotes
one crossingless circle. This differs by a factor $\delta$ from the common
normalization assigning bracket $1$ to a single circle.

The writhe $w(D)$ is computed from the incoming ports of \emph{both} strands.
If $u\in\{0,2\}$ and $o\in\{1,3\}$ are the incoming under- and over-ports,
the sign used here is $+1$ when $o-u\equiv3\pmod4$ and $-1$ otherwise.
The inherited \code{Diagram.signs()} is appropriate to its local Fox-row
convention; it is not a safe substitute for this writhe computation after
arbitrary half-turns of PD rows. That function is left unchanged to avoid
altering the Alexander backend.

A positive curl contributes $-A^3$: indeed,
$A\delta^2+A^{-1}\delta=-A^3\delta$. The other curl contributes $-A^{-3}$.
Thus the relevant unknot identity is
\begin{equation}\label{eq:unknot-bracket}
 B_D(A)=\delta(-A^3)^{w(D)}\qquad\text{when }D\text{ represents the unknot}.
\end{equation}
This is the bracket normalization underlying the Jones invariant
\cite{kauffman}. To see directly why no sign is being guessed, in the local
Temperley--Lieb algebra write $R=A1+A^{-1}e$ and
$R^{-1}=A^{-1}1+Ae$. Then
$RR^{-1}=1+(A^2+A^{-2}+\delta)e=1$, proving the R2 identity. The relations
$e_i^2=\delta e_i$ and $e_i e_{i+1}e_i=e_i$ give the R3 identity; the two
curl calculations give R1 after writhe normalization.

Evaluate \eqref{eq:bracket} at $A=2$ in $\F_p$. A discrepancy from
\eqref{eq:unknot-bracket} proves knottedness. The code never divides by
$\delta$; $A$ and $-A^3$ are nonzero, so negative exponents are well-defined.
The same routine supports any integer $A\not\equiv0\pmod p$ for testing.
There is no probability of a false positive from modular arithmetic. A
collision can only cause the filter to pass the knot on to later stages.

\subsection{Frontier dynamic programming}

At each crossing, the frontier consists of edges with exactly one processed
endpoint. A smoothing state of the processed part induces a perfect matching
of these boundary labels; every closed component is already accounted for
by a factor $\delta$. The table maps each canonical matching to the sum of
its coefficients in $\F_p$.

For each table entry and each smoothing, an independent disjoint-set
computation joins the old matching arcs to the two new smoothing arcs.
A connected component with no remaining boundary endpoints contributes a
factor $\delta$; a component with two endpoints contributes a new matching
pair. Components with any other number of endpoints signal an invariant
violation. Equal output pairings are merged by field addition and zero
coefficients are removed. Unlike the Khovanov routine, this implementation
uses no dotted-cobordism composition, no delooping matrices, and no homology
cancellation. Its independent structure is useful for testing.

\begin{proposition}
After every scan prefix, each table coefficient is exactly the sum of
partial bracket-state weights inducing that boundary matching. Therefore
the final empty-matching coefficient is \eqref{eq:bracket} modulo $p$,
independently of the scan order.
\end{proposition}
\begin{proof}
Initially the empty matching has coefficient one. Attaching a crossing
partitions full extensions uniquely by their previous state and one of
the two local smoothings. The union-find computation detects precisely
the circles closed by that attachment and the surviving pairs. Multiplying
by the local smoothing coefficient and $\delta$ for each such circle gives
the correct extended weight. Summing identical matchings is valid because
all future attachments depend only on those matchings. Induction completes
the proof; the crossingless input is handled separately as one circle.
\end{proof}

If the maximum frontier size is $w$, the number of possible pairings is at
most $(w-1)!!$, with $(-1)!!=1$. Thus a conservative uncapped bound is
\[
 \OO\!\left(n\,\operatorname{poly}(w)\,(w-1)!!\right)
 = n\,2^{\OO(w\log w)}
\]
field operations, plus scan-order construction. An arbitrary processed
region need not be a disk with all endpoints on one cyclic boundary, so
we do not silently replace this bound by a Catalan number. The default
limits are $50\,000$ live states and $1\,000\,000$ transitions. Reaching
either is an instruction to skip this filter, not to decide the knot.

The Conway and Kinoshita--Terasaka examples both have Alexander polynomial
one in the supplied data and are accepted by the first modular filter.
The new bracket test rejects both without computing their reduced
Khovanov rank $33$. Their raw residues and scan statistics are included
in the benchmark output. This is a genuine extension of the early-decision
coverage beyond Alexander-only preprocessing.

\section{Accelerating the Khovanov backend itself}
\label{sec:kh}

\subsection{The unchanged mathematical representation}

The scanner attaches one crossing at a time. Objects are boundary matchings
with a homological cube degree. A closed circle is delooped into two copies.
For matchings $a,b$, glue $a$ to the reflected $b$; each resulting circle
labels a dot variable. A morphism is a sum over $\F_2$ of square-free
monomials in those variables, represented by a set of dot-subsets. This is
the local dotted-cobordism method, with Gaussian cancellation after every
attachment \cite{bar-natan}.

We retain this representation. In particular, isomorphic matching types
are \emph{not} merged into a single chain object: their multiplicities,
cube degrees, and maps can carry essential homology. Quantum grading is
forgotten because total rank suffices for recognition, but the ordinary
chain-complex degree and differential are preserved.

\subsection{Units are self-inverse over the square-zero algebra}

For a matching with $k$ arcs, its endomorphism algebra has the form
\[
 R_k=\F_2[x_1,\ldots,x_k]/(x_1^2,\ldots,x_k^2).
\]
The unit test used by the scanner is the presence of the constant monomial.

\begin{lemma}\label{lem:unit}
Every element $u\in R_k$ with constant term $1$ satisfies $u^2=1$.
\end{lemma}
\begin{proof}
Write $u=1+v$, with $v$ a sum of nonconstant square-free monomials.
In characteristic two, all cross terms in $v^2$ vanish. Every remaining
monomial square contains some $x_i^2$ and hence vanishes. Therefore
$(1+v)^2=1$.
\end{proof}

The inverse routine consequently returns the same morphism instead of
forming a nilpotent geometric series. The tests exhaust all $128$ units
of $R_3$, using both the old and new multiplication. This shortcut depends
on the coefficient field and relations; it must not be copied unchanged
to integral or deformed Khovanov theories.

Multiplication in an endomorphism ring is also specialized: multiply two
monomials only when their dot sets are disjoint, take their union, and
XOR repeated outputs. Exact identity morphisms use an even shorter path.
General source/intermediate/target matching triples still use the complete
surface-composition routine.

\subsection{Compile topology separately from coefficients}

For $f:a\to b$ and $g:b\to c$, the connected surface components, their
Euler characteristics, output boundary circles, and assignment of input
dot-circles to components depend on $(a,b,c)$, not on the coefficients in
$f$ and $g$. The old implementation reconstructs this topology on every
composition. The new \code{\_composition\_plan} caches it, after which each
pair of monomials supplies only dot counts to the existing expansion rules.
Repeated images of differential entries under the same crossing attachment
are cached separately by \code{\_crossing\_entries}.

Four LRU caches have finite capacities of $16\,384$ entries each. Cached
coefficient collections are immutable; the public crossing-entry wrapper
returns mutable copies where mutation is needed. Caches are cleared at
entry and in a \code{finally} block at exit from each nonempty rank
calculation. This removes reliance on warm caches and prevents unbounded
retention across unrelated jobs. Entry counts are not byte limits: an
individual morphism or topological key can still be large. Internal cached
geometry is treated as read-only by the engine.

\subsection{Sparse pivot selection and Schur complements}

Suppose a differential entry $\phi:b\to c$ is an isomorphism. For any other
source $a\to c$ and target $b\to d$, cancellation replaces the corresponding
map by
\begin{equation}\label{eq:schur}
 d'_{a,d}=d_{a,d}+d_{b,d}\,\phi^{-1}\,d_{a,c}.
\end{equation}
The plus sign is subtraction over $\F_2$. In block form this is the Schur
complement, obtained by invertible row and column operations followed by
removing the contractible pair. Consequently pivot order affects work and
fill-in but not homology.

The baseline takes available unit entries from a last-in-first-out list.
The new engine uses the approximate Markowitz cost
\[
  (\#\operatorname{incoming}(c)-1)
  (\#\operatorname{outgoing}(b)-1),
\]
namely the number of possible correction pairs in \eqref{eq:schur}. A heap
favors sparse pivots, especially zero-correction pivots. Entries are checked
again when popped; stale or no-longer-invertible entries are discarded,
and changed costs are reinserted. This is a \emph{lazy heuristic}, not a
claim that every chosen pivot globally minimizes the current fill cost.
Stale keys of other heap entries may delay a better candidate.

The implementation factors $\phi^{-1}d_{a,c}$ once for each incoming source
and reuses it for all outgoing targets. Lemma~\ref{lem:unit} removes inverse
expansion. If either side is empty, no correction compositions are needed.
The extra heap overhead can lose on easy cases; the benchmark measures
rather than assumes a net gain.

The \code{compositions} statistic now counts actual composition calls.
The baseline used a different counter convention, so cross-version ratios
of that field would be misleading. Wall-clock time, computed ranks,
maximum live objects, and scan boundaries have directly comparable meanings.
The new \code{correction\_pairs} and cache-hit fields are diagnostic only.

\section{Incremental scan-order construction}

Let $s_i$ be the number of edges from a remaining crossing $i$ into the
processed region, and let $\ell_i$ be its number of self-loops. At a fixed
step with boundary size $B$, the old greedy score is
\[
 \bigl(s_i,-(B+4-2s_i-2\ell_i),-i\bigr).
\]
Lexicographically, this is equivalent to $(s_i,\ell_i,-i)$. The common
$B$ need not be recomputed for every candidate. A heap maintains the
negated score; when a crossing is added, only its remaining neighbors
need updates. Each nonloop edge raises a remaining neighbor's shared-edge
count once. Old heap records are ignored when their shared count is stale.

\begin{proposition}
For a fixed starting crossing, the new order is exactly the baseline's
greedy order and takes $\OO(n\log n)$ time and $\OO(n)$ space instead of
quadratic candidate rescanning.
\end{proposition}
\begin{proof}
Score equivalence gives identical tie breaking. The crossing graph has
bounded degree, so the total number of initial and updated heap records
is $\OO(n)$. Each extraction or insertion takes $\OO(\log n)$ time.
\end{proof}

The production calls preserve the baseline's approximate sampling policy
with \code{tries=12}. This is a \emph{target}, not a strict maximum number
of starts: the inherited integer-step rule can sample as many as $23$.
It is nevertheless a fixed constant, so choosing the best sampled order
still takes $\OO(n\log n)$. Calling \code{best\_scan\_order(tries=None)}
tries all $n$ starts and instead takes $\OO(n^2\log n)$. The implementation
and tests retain that public behavior.

Explicit orders now must be permutations of \emph{all} crossing indices.
The baseline did not enforce that condition; an empty supplied order on a
nonempty input could bypass the entire computation. Duplicate, omitted,
out-of-range, and boolean indices are rejected. This validation is important
when experimenting with alternative order heuristics.

\section{Complexity: the improvement and the remaining exponential}

\begin{center}
\begin{tabularx}{\linewidth}{@{}lXX@{}}
\toprule
Stage & Supplied implementation & New implementation\\
\midrule
R1/R2 cleanup & Repeated full searches and rebuilds; quadratic lower bound on curls
 & $\OO(n\log n)$ word operations\\
Descending test & Quadratic worst-case traversal search & $\OO(n)$ word operations\\
One greedy order & $\OO(n^2)$ candidate work & $\OO(n\log n)$ word operations\\
Determinant prefilter & Not a separate early stage & $\OO(n^3)$ fixed-field operations\\
Jones prefilter & Absent & $n2^{\OO(w\log w)}$ uncapped; bounded by work caps in production\\
Complete Kh fallback & Exponential & Exponential; less repeated algebra and fill\\
\bottomrule
\end{tabularx}
\end{center}

The $2^{\OO(n)}$ envelope for the complete fallback is deliberately coarse.
There are $2^n$ smoothing states, each with $\OO(n)$ circles and hence at
most $2^{\OO(n)}$ enhanced states. Finite matrix operations on exponentially
many generators remain within an exponential envelope over a fixed field.
Polynomial preprocessing and bounded optional filters do not change that
worst-case classification. No smaller exponential base is proved here.

\paragraph{Boundary width does not bound homology multiplicity.}
Counting boundary matchings alone is not a running-time proof for the
Khovanov scanner. Consider a connected sum of $m$ trefoils, arranged as a
sequence of fixed-size tangle blocks. A left-to-right scan has constant
frontier width. Over a field, the reduced connected-sum complex is the
tensor product of the reduced complexes, so the reduced ranks multiply.
The reduced trefoil rank is $3$, and the sum has rank $3^m$ at $n=3m$
crossings. An explicit minimal complex consequently has exponentially
many generators despite constant width. The tensor-product assertion can
also be seen directly by marking the connected-sum arc: after fixing its
label in the reduced complex, smoothing choices and the differential split
between the two factors; field-valued K\"unneth then multiplies ranks.

This example is about the explicit homology representation, not a lower
bound for every unknot decision algorithm. Our modular determinant rejects
that family cheaply. It explains why algebraic caching, width reduction,
and geometric separator intuition cannot on their own establish a
quasi-polynomial guarantee for the current fallback.

\section{Validation and reproducible measurements}
\label{sec:validation}

\subsection{Correctness checks actually executed}

The original suite passes all $19$ tests against the unchanged baseline.
The distributed suite has \TestCount\ passing tests. The legacy tests are
retained; their pipeline-method assertions explicitly disable the two new
filters, and the old resource-limit CLI check adds \code{--no-filters} so
it still reaches the stage it is intended to test. This compatibility
adjustment is documented rather than presented as an unmodified test suite.

The added checks include: exact scan-order comparisons on $100$ generated
knots and every possible starting crossing; $250$ comparisons of complete
simplification traces with independent replay; $300$ comparisons of the
linear descending test with the old brute-force search; and $400$ morphism
composition comparisons on planar boundary matchings. All $128$ units in
the three-variable endomorphism algebra are tested against the old algebra.

The modular bracket is checked against an independent full state sum for
$100$ random knots at two evaluation points and arbitrary scan orders.
The state-sum reference uses graph traversal rather than the frontier
union-find routine. Other checks compare $120$ determinant residues with
the exact Alexander polynomial, test half-turns and mirrors on $80$ inputs,
check normalized bracket invariance under available R1/R2 moves on $80$
inputs, and test $60$ guaranteed unknots generated by braid conjugation
and Markov stabilization. Mocked inconclusive residues verify that filter
collisions cannot create a false unknot verdict.

For homology, every one-component three-strand braid word of lengths $2$
and $4$ over $\{\sigma_1^{\pm1},\sigma_2^{\pm1}\}$ is checked: exactly
$168$ words. Their reduced ranks are compared with a separate dense reduced
cube implementation. Another $100$ randomly generated knots up to nine
crossings receive the same comparison. The inherited tests include further
small cube checks. Intermediate $d^2=0$ checks are enabled in the new
exhaustive and random rank tests. Tests also cover malformed orders,
invalid resource options, budget propagation, and cache cleanup.

\paragraph{Hard-case cross-check.}
The automatic greedy order and the natural input order both give reduced rank $2949$ and identical unreduced ranks in every cube degree, with intermediate $d^2=0$ checks. The respective elapsed times in these debug runs are 2.89 and 19.36 seconds. Two exploratory alternative-order checks (reverse natural order and a greedy order starting at crossing $35$) did not finish within the external run limits. No agreement is claimed for them. This illustrates important order sensitivity; these debug runs are not the three-repetition benchmark. The completed outputs are in \code{results/hard-case-verification.json}; the unsuccessful attempts are described in \code{results/additional-order-notes.txt}.

These are substantial regression checks, not a formal verification of the
implementation or an exhaustive validation at $36$ crossings. In particular,
agreement between scan orders still shares the same cobordism algebra;
the independent dense-cube checks cover small inputs, not this entire
large case. Random tests use fixed seeds recorded in the source.

\subsection{Methodology}

The comparison uses the same interpreter and host for both packages:
\PythonVersion\ on \HostPlatform. The reported processor is
\CpuModel. Each engine/input pair runs in a separate process, engines
run sequentially, and their order alternates by case. Each successful
measurement is the median of three calls with explicit algebra-cache
clearing and garbage collection before the timed region. Imports and input
construction are excluded. Raw samples and process peak RSS, when available,
are retained in \code{results/benchmark.json}. RSS is a process high-water
mark, not a per-call allocation measurement.

The corpus contains \emph{all} $25$ raw scan inputs and all $10$ pipeline
examples from the supplied benchmark, plus $10$ larger preprocessing-only
measurements. Every generated braid word is frozen in
\code{tools/corpus.json}. Raw scan measurements bypass reductions,
descending tests, and both old and new invariant filters. Pipeline
measurements use the defaults of each version. Therefore a pipeline
speedup and a raw homology speedup are not conflated.

Each call has a \TimeLimit-second cooperative limit and a $500\,000$
Khovanov-object ceiling. A parent-process timeout backs up cooperative
limits. A timed-out run is not assigned a fictitious completion time or a
median of three successful runs: its first limit event is recorded and
repetitions stop. Reported speedup lower bounds use the time limit divided
by the optimized median. The original archive's Windows/Python timing of
$600$ seconds for the hard case is preserved as historical data, but is
\emph{not} used to compute a same-machine speedup.

\subsection{Raw Khovanov results}

\small
\begin{longtable}{@{}lrrrrr@{}}
\toprule
Input family & $n$ & Reduced rank & Old (s) & New (s) & Speedup\\
\midrule
\endfirsthead
\toprule
Input family & $n$ & Reduced rank & Old (s) & New (s) & Speedup\\
\midrule
\endhead
\bottomrule\endfoot
Stabilized unknot & 8 & 1 & 0.0009 & 0.0009 & 1.03\\
Stabilized unknot & 16 & 1 & 0.0029 & 0.0023 & 1.27\\
Stabilized unknot & 32 & 1 & 0.0083 & 0.0036 & 2.32\\
Stabilized unknot & 64 & 1 & 0.0223 & 0.0065 & 3.42\\
Stabilized unknot & 128 & 1 & 0.0833 & 0.0156 & 5.34\\
Stabilized unknot & 256 & 1 & 0.3008 & 0.0265 & 11.35\\
$T(2,k)$ & 5 & 5 & 0.0036 & 0.0016 & 2.27\\
$T(2,k)$ & 11 & 11 & 0.0142 & 0.0052 & 2.72\\
$T(2,k)$ & 21 & 21 & 0.0429 & 0.0151 & 2.85\\
$T(2,k)$ & 41 & 41 & 0.2163 & 0.0336 & 6.44\\
$T(3,k)$ & 8 & 5 & 0.0122 & 0.0082 & 1.48\\
$T(3,k)$ & 10 & 7 & 0.0133 & 0.0184 & 0.72\\
$T(3,k)$ & 14 & 9 & 0.0384 & 0.0261 & 1.47\\
$T(3,k)$ & 16 & 11 & 0.0374 & 0.0314 & 1.19\\
$T(3,k)$ & 20 & 13 & 0.0653 & 0.0617 & 1.06\\
$T(3,k)$ & 22 & 15 & 0.1708 & 0.0637 & 2.68\\
random 3-braid & 20 & 75 & 0.2116 & 0.0719 & 2.94\\
random 3-braid & 40 & 321 & 0.9623 & 0.2141 & 4.49\\
random 4-braid & 15 & 5 & 0.0186 & 0.0089 & 2.08\\
random 4-braid & 21 & 25 & 0.7309 & 0.1389 & 5.26\\
random 4-braid & 31 & 45 & 0.3965 & 0.1255 & 3.16\\
random 4-braid & 41 & 109 & 1.663 & 0.4447 & 3.74\\
random 5-braid & 24 & 13 & 0.0670 & 0.0435 & 1.54\\
random 5-braid & 36 & 2949 & $>30$ & 2.233 & $>13.43$\\
random 6-braid & 31 & 85 & 1.008 & 0.1245 & 8.10\\
\end{longtable}
\normalsize

All successfully completed baseline/optimized comparisons agree on rank.
The hard case has \HardRank\ reduced generators and \HardUnreduced\
unreduced generators. Its optimized maximum object count before elimination
is \HardPeakObjects, and its maximum frontier contains \HardWidth\ edges.
These counts emphasize that the improvement does not arise from secretly
replacing the hard diagram by a lower-crossing representative in the raw
benchmark. The scan order is the baseline's own selected order.

\subsection{Full pipeline results}

\begin{center}\small
\begin{tabular}{@{}lrrrl@{}}
\toprule
Example & Old (ms) & New (ms) & Speedup & New decisive stage\\
\midrule
Conway & 47.693 & 1.622 & 29.41 & Bracket\\
Figure eight & 0.137 & 0.162 & 0.85 & Determinant\\
Determinant-one grid & 0.759 & 0.807 & 0.94 & Bracket\\
Scrambled grid unknot & 0.376 & 0.234 & 1.61 & R1/R2\\
Hard unknot (8) & 2.699 & 2.146 & 1.26 & Khovanov\\
Kinoshita--Terasaka & 48.491 & 1.004 & 48.28 & Bracket\\
$T(3,5)$ & 0.706 & 1.006 & 0.70 & Bracket\\
Trefoil & 0.117 & 0.153 & 0.76 & Determinant\\
Unknot example & 0.023 & 0.073 & 0.31 & R1/R2\\
Unknot braid (40) & 7.050 & 0.404 & 17.47 & R1/R2\\
\bottomrule\end{tabular}\end{center}

The measurements were made on a shared container host; scheduling noise can
also produce non-monotone individual timings. No asymptotic exponent is fitted
to these samples. The easiest examples may be slower because the new bookkeeping, input digest,
and modular-stage dispatch have fixed overhead. Sub-millisecond differences
should not be overinterpreted. The determinant-one examples are more
informative: they demonstrate that the bracket stage can avoid the complete
homology calculation on cases where the Alexander stage cannot decide.
The supplied grid example is also retained rather than choosing only named
low-crossing knots.

\subsection{Preprocessing scaling}

\begin{center}\small
\begin{tabular}{@{}lrrrr@{}}
\toprule
Subroutine & $n$ & Old (s) & New (s) & Speedup\\
\midrule
Scan ordering & 128 & 0.0700 & 0.0026 & 26.6\\
R1/R2 cleanup & 128 & 0.0733 & 0.0009 & 86.0\\
Scan ordering & 256 & 0.2694 & 0.0052 & 52.1\\
R1/R2 cleanup & 256 & 0.7270 & 0.0016 & 446.9\\
Scan ordering & 512 & 1.088 & 0.0248 & 43.9\\
R1/R2 cleanup & 512 & 1.107 & 0.0032 & 343.1\\
Scan ordering & 1024 & 4.201 & 0.0211 & 198.8\\
R1/R2 cleanup & 1024 & 4.883 & 0.0078 & 623.8\\
Scan ordering & 2048 & 20.202 & 0.0531 & 380.3\\
R1/R2 cleanup & 2048 & 25.057 & 0.0138 & 1818.4\\
\bottomrule\end{tabular}\end{center}

These runs use the closure of
$\sigma_1\sigma_2\cdots\sigma_n$ on $n+1$ strands. For order construction,
only \code{best\_scan\_order(tries=12)} is timed. For reduction, only the
simplifier is timed, and both versions remove every crossing. The family
has $n$ crossings and linear PD size; braid construction is outside the
timed region. The measurements illustrate the proved subroutine bounds,
not a new asymptotic theorem about unknot recognition.

\clearpage
\section{Using and extending the package}

From the extracted archive root, Python $3.10$ or later is sufficient.
There are no runtime packages to install:
\begin{lstlisting}
python -m fastunknot recognize examples/conway.json
python -m fastunknot khovanov examples/conway.json --check-d2
python -m unittest discover -s tests -v
python tools/benchmark.py --limit 30 --repeats 3
python tools/verify_hard_case.py
\end{lstlisting}
The installed \code{fastunknot} console command is also available after
\code{python -m pip install .}. For a direct Python call:
\begin{lstlisting}[language=Python]
from fastunknot import Diagram, recognize

diagram = Diagram.from_braid(2, [1, 1, 1])
result = recognize(diagram, seconds=30, max_objects=100_000)
print(result.status, result.method)
\end{lstlisting}

\code{--no-filters} disables both new modular stages while keeping the
accelerated simplifier and Khovanov engine. \code{--no-alexander} disables
only the full symbolic Alexander calculation; it does not disable the
modular determinant. \code{--no-reduction} and \code{--no-descending}
control the old positive-decision shortcuts. The \code{khovanov} subcommand
bypasses every preprocessing stage. The \code{jones} subcommand reports a
modular comparison, not a complete unknot verdict. Exit code $0$ means a
completed exact recognition result (either kind), $2$ means invalid input
or options, and $3$ means a resource limit. Inspect the JSON status rather
than treating every nonzero return as knottedness.

Time limits are cooperative. A large Python operation, input validation,
order construction, or the inherited exact Alexander calculation can
overshoot a requested deadline. The object ceiling counts chain objects,
not bytes, dictionary entries, morphism terms, or cache payload. Untrusted
large jobs should run in a separately supervised process with operating-
system memory and wall-time limits. Cache clearing is safe for algebraic
correctness but concurrent calls sharing a process can disrupt each
other's cache-hit statistics and performance; the benchmark uses isolated
sequential processes.

The highest-value next optimization experiment would be compact integer
bitsets for dot monomials and interned local matching identifiers. It must
preserve the coefficient/topology separation and the multiplicity of chain
objects. Componentwise or connected-sum decomposition could avoid materializing
some product homologies, but such a change needs a correct decomposition
certificate and must not be inferred from repeated matching types alone.
Neither direction is included in the claimed measurements or complexity
improvements of this release.

\section{Conclusion}

The requested quasi-polynomial target is not established by this code.
The fallback goal is achieved in two independent ways: several preprocessing
subroutines have better proved asymptotic costs, and the complete supplied
homology engine is substantially faster on the measured corpus. A previously
unfinished $36$-crossing homology computation now completes, while a bounded
bracket stage provides fast certified rejection of determinant-one examples.
The decisive distinction is between reducing repeated work and fill-in,
which is implemented and measured here, and compressing the geometric
recognition problem enough to prove $n^{\OO(\log n)}$, which requires a
different set of invariants and algorithmic bounds.

\begin{thebibliography}{9}
\bibitem{lackenby-talk}
M. Lackenby, \emph{Unknot recognition in quasi-polynomial time}, talk slides,
main theorem labeled 2021, 109 pages. Supplied as
\code{docs/quasipolynomial-talk.pdf}. Online copy accessed 18 September 2026.
\url{https://people.maths.ox.ac.uk/lackenby/quasipolynomial-talk.pdf}.

\bibitem{lackenby-paper}
M. Lackenby, \emph{Incompressible surfaces, hierarchies and unknot recognition},
arXiv:2607.23350v1, 25 July 2026. Sections 9--10 and introduction.
The supplied source is in \code{docs/arXiv-2607.23350v1/}.
\url{https://arxiv.org/abs/2607.23350}.

\bibitem{bar-natan}
D. Bar-Natan, \emph{Fast Khovanov homology computations},
Journal of Knot Theory and Its Ramifications \textbf{16} (2007), 243--255.
arXiv:math/0606318; delooping and Gaussian elimination are Lemmas 4.1 and
4.2 in the arXiv version.
\url{https://arxiv.org/abs/math/0606318}.

\bibitem{kronheimer-mrowka}
P. B. Kronheimer and T. S. Mrowka, \emph{Khovanov homology is an
unknot-detector}, Publications Math\'ematiques de l'IH\'ES \textbf{113}
(2011), 97--208. \url{https://arxiv.org/abs/1005.4346}.

\bibitem{shumakovitch}
A. N. Shumakovitch, \emph{Torsion of the Khovanov homology},
arXiv:math/0405474v2, 2018. Corollary 3.2.C.
\url{https://arxiv.org/abs/math/0405474}.

\bibitem{kauffman}
L. H. Kauffman, \emph{State models and the Jones polynomial},
Topology \textbf{26} (1987), 395--407.
\url{https://doi.org/10.1016/0040-9383(87)90009-7}.
\end{thebibliography}
\end{document}
